% Fichier généré par rapport/benchmark_akori.py : ne pas modifier à la main.
\begin{longtable}{@{}r p{0.50\textwidth} c c c c@{}}
\caption{Résultat de chaque question du banc de test.}\label{tab:questions}\\
\toprule \# & Question & Sim. max & Passage & Seuil & Réponse \\ \midrule \endfirsthead
\toprule \# & Question & Sim. max & Passage & Seuil & Réponse \\ \midrule \endhead
\bottomrule \endfoot
\multicolumn{6}{@{}l}{\textbf{Dans le cours}} \\
1 & Quel modèle d'embedding utilise AKORI ? & 0{,}68 & oui & accepté & \nonmesure \\
2 & Quelle bibliothèque assure la recherche vectorielle des passages ? & 0{,}57 & oui & accepté & \nonmesure \\
3 & Quelle est la taille des fragments lors du découpage du texte ? & 0{,}70 & oui & accepté & \nonmesure \\
4 & Quel chevauchement est conservé entre deux fragments successifs ? & 0{,}65 & oui & accepté & \nonmesure \\
5 & Combien de passages sont transmis au modèle génératif pour une question ? & 0{,}63 & oui & accepté & \nonmesure \\
6 & Quel est le seuil de pertinence utilisé par défaut ? & 0{,}57 & oui & accepté & \nonmesure \\
7 & Quel modèle génératif principal est utilisé par AKORI ? & 0{,}63 & non & accepté & \nonmesure \\
8 & Quelle méthode protège les mots de passe des comptes ? & 0{,}60 & oui & accepté & \nonmesure \\
9 & Combien d'itérations utilise la dérivation des mots de passe ? & 0{,}57 & oui & accepté & \nonmesure \\
10 & Quelle bibliothèque extrait le texte des fichiers PDF ? & 0{,}57 & oui & accepté & \nonmesure \\
11 & Quel moteur OCR lit les pages scannées ? & 0{,}59 & oui & accepté & \nonmesure \\
12 & Avec quel outil l'interface web a-t-elle été construite ? & 0{,}56 & oui & accepté & \nonmesure \\
13 & Quel modèle de langage local a été testé au début du projet ? & 0{,}54 & oui & accepté & \nonmesure \\
14 & Quelle latence observait le prototype exécuté localement ? & 0{,}61 & oui & accepté & \nonmesure \\
15 & Quels sont les intervalles de la répétition espacée du plan de révision ? & 0{,}64 & oui & accepté & \nonmesure \\
16 & Quelle est la dimension des vecteurs produits par BGE-M3 ? & 0{,}66 & oui & accepté & \nonmesure \\
17 & Que signifie le sigle RAG ? & 0{,}54 & oui & accepté & \nonmesure \\
18 & Sur combien de points la colle orale est-elle notée ? & 0{,}59 & oui & accepté & \nonmesure \\
19 & Où sont enregistrées les données lorsque l'application tourne sur Colab ? & 0{,}59 & oui & accepté & \nonmesure \\
\midrule
\multicolumn{6}{@{}l}{\textbf{Reformulées}} \\
20 & Quel outil compare les vecteurs pour retrouver les extraits les plus proches ? & 0{,}58 & oui & accepté & \nonmesure \\
21 & Avec quel modèle les textes sont-ils transformés en listes de nombres ? & 0{,}57 & oui & accepté & \nonmesure \\
22 & Comment l'application lit-elle les pages qui ne sont que des images ? & 0{,}64 & oui & accepté & \nonmesure \\
23 & Que fait l'assistant quand aucun extrait n'est assez proche de la demande ? & 0{,}55 & non & accepté & \nonmesure \\
24 & Pourquoi a-t-on renoncé à faire tourner le modèle de langage sur la machine ? & 0{,}49 & oui & accepté & \nonmesure \\
25 & Comment les secrets de connexion des utilisateurs sont-ils conservés de façon sûre ? & 0{,}58 & oui & accepté & \nonmesure \\
26 & Quelle est la longueur de chaque morceau de texte indexé ? & 0{,}59 & oui & accepté & \nonmesure \\
27 & Sur quel principe de psychologie reposent les cartes mémoire et les questionnaires ? & 0{,}51 & non & accepté & \nonmesure \\
28 & Comment l'application continue-t-elle de répondre quand le service d'IA est saturé ? & 0{,}54 & non & accepté & \nonmesure \\
29 & Où peut-on héberger l'application pour qu'elle ait une adresse permanente ? & 0{,}56 & oui & accepté & \nonmesure \\
\midrule
\multicolumn{6}{@{}l}{\textbf{Proches du sujet, réponse absente du cours}} \\
30 & Quelle est la taille maximale d'un PDF importé ? & 0{,}62 & -- & accepté & \nonmesure \\
31 & Combien coûte un abonnement mensuel à AKORI ? & 0{,}59 & -- & accepté & \nonmesure \\
32 & Quelle version de Python est utilisée dans l'image Docker ? & 0{,}55 & -- & accepté & \nonmesure \\
33 & Sur quel port réseau l'application écoute-t-elle sur le serveur ? & 0{,}53 & -- & accepté & \nonmesure \\
34 & Combien d'étudiants ont testé AKORI pendant le semestre ? & 0{,}62 & -- & accepté & \nonmesure \\
35 & AKORI fonctionne-t-il sur un Raspberry Pi ? & 0{,}64 & -- & accepté & \nonmesure \\
\midrule
\multicolumn{6}{@{}l}{\textbf{Hors sujet}} \\
36 & Quelle est la capitale de l'Australie ? & 0{,}30 & -- & refus & \nonmesure \\
37 & Comment fonctionne la photosynthèse chez les plantes ? & 0{,}43 & -- & accepté & \nonmesure \\
38 & Qui a peint le tableau de la Joconde ? & 0{,}28 & -- & refus & \nonmesure \\
39 & Quelle est la formule de l'aire d'un cercle ? & 0{,}39 & -- & accepté & \nonmesure \\
40 & Expliquez la théorie de la relativité restreinte d'Einstein. & 0{,}44 & -- & accepté & \nonmesure \\
41 & Comment préparer une pâte à crêpes ? & 0{,}41 & -- & accepté & \nonmesure \\
42 & Quelles sont les causes de la Première Guerre mondiale ? & 0{,}36 & -- & accepté & \nonmesure \\
43 & Quel est le rôle des mitochondries dans la cellule ? & 0{,}41 & -- & accepté & \nonmesure \\
44 & Comment calculer la dérivée d'une fonction polynomiale ? & 0{,}41 & -- & accepté & \nonmesure \\
45 & Qu'est-ce que le protocole de routage OSPF ? & 0{,}44 & -- & accepté & \nonmesure \\
46 & Quel est le plus long fleuve d'Afrique ? & 0{,}33 & -- & accepté & \nonmesure \\
47 & Comment fonctionne un moteur à explosion ? & 0{,}45 & -- & accepté & \nonmesure \\
48 & Quelles sont les règles du hors-jeu au football ? & 0{,}37 & -- & accepté & \nonmesure \\
49 & Qu'est-ce que l'inflation en économie ? & 0{,}42 & -- & accepté & \nonmesure \\
50 & Comment soigner une grippe saisonnière ? & 0{,}35 & -- & accepté & \nonmesure \\
51 & Quelle est la vitesse de la lumière dans le vide ? & 0{,}41 & -- & accepté & \nonmesure \\
52 & Qui a écrit le roman Les Misérables ? & 0{,}30 & -- & refus & \nonmesure \\
53 & Comment fonctionne la blockchain du Bitcoin ? & 0{,}46 & -- & accepté & \nonmesure \\
54 & Quelles sont les étapes de la mitose ? & 0{,}42 & -- & accepté & \nonmesure \\
55 & Comment se forme un arc-en-ciel ? & 0{,}41 & -- & accepté & \nonmesure \\
\end{longtable}
